\input{./._relpath-to-latexroot.ltx}
\documentclass[\latexroot/\projectname]{subfiles}
\input{\latexroot/@resources/subfile-setup.ltx}

\begin{document}

\begin{table}[h]
  \centering
  \caption{Local auto spending multipliers\whenintegrated{\label{tab:empirics-auto-multipliers}}}
  \begin{adjustbox}{width=\textwidth}
  \begin{tabular}{lllll}
    \hline
     & \multicolumn{4}{c}{Auto spending} \\
     & \multicolumn{4}{c}{(CCP/Equifax)} \\
    \cline{2-5}
    \textbf{Spending category} & \textbf{OLS} & \textbf{IV} & \textbf{OLS} & \textbf{IV} \\
    \hline
    Government spending & 0.07*** (0.01) & 0.11*** (0.01) & 0.06*** (0.01) & 0.09*** (0.01) \\
    Partial $F$ stat. & - & 248.5 & - & 187.9 \\
    County controls/state fixed effects & No & No & Yes & Yes \\
    \# of counties & 3,119 & 3,119 & 3,047 & 3,047 \\
  \end{tabular}
  \end{adjustbox}
  \begin{flushleft}
    \footnotesize
    \textit{Notes:} The table shows the estimates of the regression of the growth rate in auto spending on cumulative government spending at the county level during the period 2008-12. The standard errors are given in parentheses and are clustered at the state level. One, two, and three stars denote significance at the $10 \%, 5 \%$, and $1 \%$ levels, respectively.
  \end{flushleft}
\end{table}

\smartbib

\end{document}
